Diabetic retinopathy is a complication of diabetes that may progress silently and cause irreversible vision damage when not identified and treated in time. It is one of the leading causes of avoidable blindness among people with diabetes in Brazil, which makes early diagnosis essential. Through fundus examination, relevant indicators can be identified, such as soft and hard exudates, which help assess the disease. However, in several regions access to specialists and adequate infrastructure for early diagnosis remains limited, in addition to the subjectivity involved in interpreting the examination, which varies even among specialists and is greater when it is performed by non-specialist professionals. This work consists of a comparative study of different deep learning architectures, using traditional models, such as VGG16 and ResNet50, lightweight models, such as MobileNetV3 and EfficientNet-B0, and Transformer-based models, such as ViT and Swin Transformer, in order to identify the most suitable models for the segmentation of hard and soft exudates and for the classification of the diabetic retinopathy grade. In segmentation, the six networks will be used as encoders of a single U-Net, trained on IDRiD and also tested on DDR, and in classification they will be trained to grade the disease from 0 to 4 using the APTOS and EyePACS datasets, and also tested on Messidor-2 without any fine-tuning. All models will be trained under the same conditions, with cross-validation, and compared by statistical tests with correction for multiple comparisons and by computational cost, measured by the number of parameters and operations, inference time and memory. The selected models will be used to build a computer-aided diagnosis (CAD) solution that shows the lesions, the predicted grade and a macular edema alert, in order to support the accurate and efficient diagnosis of diabetic retinopathy in regions with limited medical resources.

% Separe as Keywords por ponto e vírgula.
\keywords{diabetic retinopathy; image segmentation and classification; deep learning; exudates; computer-aided diagnosis.}
